\documentclass[11pt]{article}
\usepackage{amsmath,amssymb,amsthm,hyperref}
\newtheorem{proposition}{Proposition}
\title{Symmetry, Manin classes, and a norm-one period countercontrol}
\author{Internal literature-bridge draft}
\date{16 November 2026}
\begin{document}\maketitle
These deductions are unregistered and depend on the explicitly named
published or retained inputs. They make no historical priority claim.
Neither full HC4/JC2 nor either full exceptional MZ problem is closed.
External review remains deferred; checking these deductions is internal work.

\begin{proposition}[Centered even HC4 potentials; conditional on HCTANG2]
Assume the retained Tang cotangent alternative HCTANG2.
Let $h\in\mathbb C[X_1,\ldots,X_4]$ have nonzero constant Hessian determinant.
If there are $a,b\in\mathbb C^4$ such that
$q(Z)=h(a+Z)-b\cdot Z$ satisfies $q(-Z)=q(Z)$, then $\nabla h$ is a
polynomial automorphism. No degree bound is imposed.
\end{proposition}
\begin{proof}
Translations and subtracting a linear form preserve the Hessian and
invertibility of the gradient. Suppose $\nabla q$ is not invertible.
HCTANG2 gives a constant invertible linear frame $(z,w)$ in which
$q=D(z)+w\cdot K(z)$, where $z,w\in\mathbb C^2$ and
$\det DK\in\mathbb C^*$. Central parity commutes with this linear frame.
Comparing coefficients in
$D(-z)-w\cdot K(-z)=D(z)+w\cdot K(z)$ shows that $K$ is odd.
Miyanishi's published equivariant Jacobian theorem, applied to the small
group $\{I,-I\}\subset\mathrm{GL}_2$, makes this plane Keller map an
automorphism. Then $\nabla q$ is invertible: from its outputs
$K(z)$ and $\nabla D(z)+DK(z)^t w$, recover $z$ by the polynomial
inverse of $K$, then $w$ by the polynomial inverse transpose of $DK$.
This contradicts the alternative assumed.
\end{proof}
Miyanishi's theorem is in Transformation Groups 28 (2023), 951--971,
\href{https://doi.org/10.1007/s00031-022-09727-7}{DOI}; the explicit odd-map
corollary is on page3 of \href{https://arxiv.org/abs/2110.06709v1}{the preprint}.
The published theorem is not a new result of this repository.
The universal four-variable deduction retains the unreviewed HCTANG2
classification dependency. For explicitly supplied cotangent potentials
with odd $K$ no classification is needed.

\begin{proposition}[A differential test in the split semiabelian case]
Let $K=\mathbb C(t)^{\rm alg}$ with its usual derivation. Suppose a
semiabelian variety is split as $G=A\times\mathbb G_m^r$ and a section
is $(s,u_1,\ldots,u_r)\in G(K)$. It has constant type up to isogeny in
the sense of the retained boundary-motive note if and only if
\[
 M_{K,A}(s)=0,\qquad \partial u_i/u_i=0\quad(1\leq i\leq r),
\]
where $M_{K,A}$ is the CLASSICAL differential-arithmetic Manin map.
\end{proposition}
\begin{proof}
The classical theorem of the kernel gives
$\ker M_{K,A}=\operatorname{Tor}(A)+A_0(\mathbb C)$, where $A_0$ is
the complex trace. The logarithmic derivative has kernel
$\mathbb C^*$ on $K^*$. Thus the displayed conditions imply that an
integer multiple of the section is the image of a constant point of
$A_0\times\mathbb G_m^r$. Conversely, a homomorphism from a constant
semiabelian variety into $A$ factors through its abelian quotient and
the complex trace, while its torus-coordinate characters take constant
points to constants. Multiplying by an isogeny integer does not change
the vanishing conditions. This proves the equivalence.
\end{proof}
The kernel statement and its stronger traceless complex-linear
injectivity are discussed immediately before and in Proposition2.15 of
\href{https://arxiv.org/abs/1502.02012}{Bertrand--Pillay}.
Their Proposition2.14 gives a counterexample for complex-linear
injectivity of the DIFFERENTIAL-ALGEBRAIC Manin map $\mu_A$.
The maps must not be identified. A nonsplit semiabelian extension
requires its extension data; this proposition does not remove it.
Exceptional purity has not been shown to annihilate the displayed
Manin class. That is a concrete new target, not an imported conclusion.

\begin{proposition}[The known norm-invisible pure failure is Pell-type]
In the retained Long composite countercontrol put
$\tau^2=r+1/4$ and $g=(\tau-1/2)/(\tau+1/2)$. Then
\[
 g=A(r)+B(r)\tau,\quad A(r)=1+\frac1{2r},\quad B(r)=-\frac1r,
 \qquad A(r)^2-(r+1/4)B(r)^2=1.
\]
Thus the constant norm is a nonconstant rational Pell-type solution,
with its nontrivial section in the norm-one torus.
\end{proposition}
\begin{proof}
Multiply the ratio defining $g$ by $(\tau-1/2)/(\tau-1/2)$ and use
$\tau^2-1/4=r$. The displayed norm identity follows by multiplication
with the conjugate $\tau\mapsto-\tau$.
\end{proof}
The retained note
\texttt{ALGEBRAIC\_BOUNDARY\_ONE\_MOTIVE\_AND\_CONSTANT\_SECTION.md}
gives the actual phase $P=F(F+1)$ and fixed marker $Q=-y$, with
$\mathcal L(QP^m)=(-1)^{m-1}(m!)^2/(2m+1)!\ne0$ at every $m\geq1$.
This is a reinterpretation of that existing countercontrol, not a new
exceptional phase or a polynomial Pell theorem. It motivates retaining
all branch characters. The related relative Manin--Mumford literature
\href{https://arxiv.org/abs/1307.1008}{includes genuinely semiabelian
Ribet exceptions}; the elementary toric example here is not a Ribet section.
\end{document}
